% !TEX root = /home/eb/.vscode/extensions/qftrules.programme-colle/tmp/Exercice.tex
\begin{exo}[PC][1][devoir]{La lunette astronomique}
	\Annale[Physique PC]{CCINP}{2015}
La lunette astronomique est un système contré constitué d'un objectif et d'un oculaire. L'objectif est assimilé à une lentille mince convergente de centre optique $O_1$, de distance focale $f_1'$ et de diamètre $D_1$. L'oculaire est une lentille mince convergente de centre optique $O_2$, de distance focale $f_2'$ et de diamètre $D_2$. L'objectif donne d'un objet éloigné une image réelle appelée image objective. Cette dernière est observée au moyen de l'oculaire.

\begin{quest}
\question À quelle condition l'œil d'un observateur, supposé sans défaut, n'accommode pas (ne se fatigue pas) ? En déduire la position relative de l'objectif et de l'oculaire. Ce système optique possède-t-il des foyers ? Comment se nomme un tel système optique ?
\sol{
    Un œil normal n'a pas besoin d'accommoder s'il observe un objet situé à l'infini.
    L'image objective doit donc se situer en $F_1'$. Or l'objet éloigné peut être considéré à l'infini donc cette image objective se trouve en $F_2'$. Ainsi $F_1'=F_2$.
   Ce système ne possède pas de foyer image (l'image d'un objet à l'infini est elle-même à l'infini) ni objet (l'objet donnant une image à l'infini est lui-même à l'infini). On parle donc de système afocal.
}

\question[2] Rappeler les conditions de Gauss. Réaliser un schéma, sans respecter les échelles, montrant le devenir d'un rayon incident faisant un angle $\theta$ avec l'axe optique et émergeant avec un angle $\theta'$ dans les conditions de Gauss.
\sol{
Conditions de Gauss : rayons peu inclinés par rapport à l'axe optique ($\theta\ll1$) et proches de l'axe optique (par rapport au rayon de la lentille, le rayon doit arriver près du centre, pas près des bords).
\begin{center}
\includegraphics[width=0.5\textwidth]{theta.png}
\end{center}
}

\newcommand{\Width}{0.7\textwidth}
\begin{fig}
	\centering
\includegraphics[width=\Width]{astro1.png}
%\caption{}
% \label{fig:}
\end{fig}
%%%%%%%%%%%%%

\question Déterminer l'expression du grossissement $G=\theta'/\theta$ en fonction de $f_1'$ et de $f_2'$ et calculer  ce grossissement si $f_1' = \SI{1.0}{m}$ et $f_2'=\SI{20}{mm}$.

\sol{
On a $\tan \theta = \overline{A'B'}/f_1'$  et $\tan \theta' = -\overline{A'B'}/f_2'$ (signes cohérents avec le cas de la figure) et donc $G=\frac{\theta'}{\theta} \simeq -\frac{f_1'}{f_2'}$. A.N. $G= -50<0$
cohérent avec l'inversion du sens de l'image.
}

\question[2] On considère un faisceau lumineux issu d'un point objet $A$ à l'infini sur l'axe optique de la lunette. Sans respect des échelles, représenter le devenir d'un tel faisceau lumineux limité par la monture de la lentille objectif (encore appelé diaphragme d'ouverture).

\sol{
\begin{center}
\includegraphics[width=0.5\textwidth]{monture.png}
\end{center}
}

\begin{fig}
\centering
\includegraphics[width=\Width]{astro2.png}
%\caption{}
% \label{fig:}
\end{fig}
%%%%%%%%%%%%%
\question Exprimer le diamètre $D$ du faisceau de rayons issu de l'oculaire en fonction du grossissement $G$ de la lunette ainsi que du diamètre $D_1$ du diaphragme d'ouverture.

\sol{
Thalès : $D_1/D = f_1'/f_2'$ donc $D=D_1/|G|$.
}

\question Après avoir calculé la valeur numérique du diamètre $D$ du faisceau  de rayons issu de l'oculaire, montrer que c'est le diaphragme d'ouverture, de diamètre $D_1$, qui le limite et non l'oculaire de diamètre $D_2$. On donne $D_1=\SI{10}{cm}$ et $D_2 = \SI{6}{mm}$.

\sol{
Limitation du faisceau émergent : si $D>D_2$, c'est l'oculaire qui limite la taille du faisceau émergent. Sinon, c'est l'objectif. A.N. $D=\SI{2}{mm}<D_2$.      La taille de l'objectif limite le diamètre du faisceau émergent.
}


\question On considère un objet ponctuel situé à l'infini en dehors de l'axe optique et  dans la direction $\theta$  par rapport à ce dernier. Expliquer, de façon qualitative, ce qu'il advient des rayons lumineux lorsque l'angle $\theta$ devient trop important. On dit de la monture de l'oculaire qu'elle est le diaphragme de champ de la lunette. Pouvez-vous justifier cette affirmation ?

\sol{
Quand le faisceau lumineux s'incline, il ne vient plus frapper l'oculaire (L2) de façon centrée. S'il est trop incliné, tout ou partie du faisceau \guill{tape} en dehors de (L2), sur la monture. Ils ne ressortent pas.
L'oculaire définit donc le diaphragme de champ, c'est-à-dire la zone (le champ) d'espace visible avec la lunette.
}


\begin{fig}
\centering
\includegraphics[width=\Width]{astro3.png}
%\caption{}
% \label{fig:}
\end{fig}
%%%%%%%%%%%%%

\question  L'objectif d'une lunette astronomique doit être capable de donner une image parfaite d'un point infiniment éloigné. Pour cela, il doit, notamment, être achromatique. D'où provient l'aberration chromatique   d'une lentille ? Comment, en physique, qualifie-t-on ce type de milieu ?
\sol{
Aberration chromatique : les rayons traversant la lentille sont déviés différemment selon leur longueur d'onde car l'indice optique en dépend (la focale dépend de l'indice ) ; l'image d'un point est donc une tache en lumière polychromatique, les différentes longueurs d'onde la constituant ne convergeant pas au même endroit.
    $n(\lambda)$  : milieu dispersif.
}

\end{quest}
\end{exo}
%%%%%%%%%%%%%%%%%%%%%%%%%%%%%%%%%%%


%%%%%%%%%%%%%%%%%%%%%%%%%%%%%%%%%%%%%%%%%%%%%%%%%%%%%
\begin{exo}[MPSI]{Lunette de Galilée}
Une lunette de Galilée comprend un objectif assimilable à une lentille mince $L_1$ de centre $O_1$ et de vergence $V_1 = \SI{5.0}{\delta}$ et un oculaire assimilable à une lentille mince $L_2$ de centre $O_2$ et de vergence $V_2 = \SI{-20}{\delta}$.


\begin{quest}
\item Déterminer la nature des deux lentilles et donner la valeur des distances focales $f'_1$ et $f'_2$.

\sol{
$f_1'=1/V_1=20\,$cm et $f_2'=1/V_2=-5\,$cm.
}

\item La lunette est de type afocale. Préciser dans ces conditions la position relative des deux lentilles en donnant la valeur de $d = \overline{O_1 O_2}$.

\sol{
%\includegraphics[width=0.75\textwidth]{1.png}\\
$d=f_1'+f_2'$.
}

\item Tracer, dans les conditions de Gauss, la marche d’un rayon lumineux incident arrivant d’un objet à l’infini, faisant un angle $\theta$ avec l’axe optique et émergeant sous l’angle $\theta'$.


\item En déduire le grossissement ou grandissement angulaire en fonction de $\theta$ et $\theta'$ puis des distances focales des deux lentilles. Donner sa valeur numérique.

\sol{
$G=\theta'/\theta=-f_1'/f_2'=12$ image à l'endroit.
}

\item Un astronome amateur utilise cette lunette normalement adaptée à la vision d’objets terrestres pour observer deux cratères lunaires : Copernic de diamètre 96 km et Clavius de
diamètre 240 km. On rappelle que la distance entre la Terre et la Lune vaut $384\cdot 10^6$ m. Voit-il ces deux cratères à l'œil nu ? à l’aide de la lunette ?

\sol{
$\theta_{\rm nu} \simeq d/D$ soit $0.85'$ et $2.14'$. La résolution de l'œil est de $1'$. Avec la lunette, on a $10.2'$ et $25.7'$.
}

%\item Déterminer le grandissement du dispositif.
%
%\sol{Lagrange-Helmholtz pour une lentille $G\gamma=1$.}

\end{quest}
\end{exo}
%%%%%%%%%%%%%%%%%%%%%%%%%%%%%%%%%%%%%%%%%%%%%%%%%%%%%

%%%%%%%%%%%%%%%%%%%%%%%%%%%%%%%%%%%%%%%%%%%%%%%%%%%%%
\begin{exo}[MPSI]{Le microscope}
Un microscope optique permet d’observer des globules sanguins.
Il est modélisable par deux lentilles minces convergentes $L_1$ pour l’objectif de distance focale
$f_1'$ et $L_2$ pour l’oculaire de distance focale $f_2'$. Il est réglé pour donner une image à l’infini d’un
objet réel AB perpendiculaire à l’axe optique, A étant sur l’axe optique) légèrement en avant
du foyer objet de l’objectif. Cette image est observée par un œil emmétrope (normal) placé au
voisinage du foyer image de l’oculaire. On notera A'B' l’image intermédiaire.
Le microscope porte les indications suivantes :
	\begin{itemize}
	\item $\times 40$ pour l’objectif, ce qui signifie que la valeur absolue du grandissement de l’objet AB par l’objectif est de 40 ;
	\item $\times 10$ pour l’oculaire, ce qui signifie que le grossissement ou rapport entre l’angle
sous lequel on voit l’image à l’infini d’un objet à travers l’oculaire seul et l’angle sous lequel
on voit ce même objet à l’œil nu lorqu’il est situé à la dimension minimale de vision distincte
$d=25\,$cm vaut 10 ;
 \item $\omega_0 = 0,\!65$ pour l’ouverture numérique ou valeur de $n \sin u$ avec $n$ l'indice de réfraction du milieu dans lequel se trouve l’objectif et $u$ l’angle maximum des rayons issus de A arrivant sur l’objectif.
\item  $\Delta = 16\,$cm pour l’intervalle optique ou distance entre le foyer image $F_1'$ de l’objectif et le
foyer objet $F_2$ de l’oculaire.
	\end{itemize}

\begin{quest}
	\item Faire un schéma du dispositif (sans respecter l’échelle) et tracer la marche de trois rayons
lumineux issus du point B de l’objet AB, l’un émis parallèlement à l’axe optique, un autre
passant par le foyer objet de l’objectif, et au autre passant par le centre de la lentille $L_1$.
\sol{
\includegraphics[width=0.75\textwidth]{trace_micro.png}
}
	\item En utilisant le grossissement de l'oculaire, déterminer la distance focale $f_2'$ de l’oculaire.

	\sol{
	$G=d/f_2'$ d'où $f_2'=d/G=2,\!5\,$cm
	}

	\item Déterminer la distance focale $f_1'$ de l’objectif. On pourra utiliser le grandissement $\gamma_1$ de l’objectif.

\sol{
%\includegraphics[scale=0.2]{g1.png}
 et $F_1'F_2 \cdot F_1A = - f_1'^2$ d'où $\gamma_1 = -\Delta/f_1'$ (ou même directement à partir du rayon parallèle). Donc $f_1'=-\Delta/\gamma_1=4\,$mm.
}


\item En s'aidant de la formule de conjugaison de Newton, calculer la distance $\overline{O_1A}$ permettant de positionner l’objet.

\sol{
%\includegraphics[scale=0.2]{O1A.png}
 soit $F_1 A =-f_1'^2/\Delta=-0.1\,$mm et $O_1 A=-4.1\,$mm.
}

\item Calculer le grossissement commercial pour une image finale à l’infini.
\sol{
%\includegraphics[scale=0.2]{a.png}\includegraphics[scale=0.2]{ap.png}
%\includegraphics[scale=0.2]{AF2.png}\includegraphics[scale=0.2]{G.png}
}
\item Déterminer la latitude de mise au point, à savoir la variation de la distance $\overline{O_1A}$ permettant de passer d'une observation à l'infini (ce qui a été supposé jusqu'à présent) à une observation à la distance minimale que l'œil peut accomoder, à savoir $d=25\,$cm. L'image A'B' se trouve alors à une distance $d$ à gauche de l'oculaire (elle est virtuelle).

\item Calculer l’angle u intervenant dans l’ouverture numérique pour un objectif placé dans
l’air. Le microscope est-il utilisé dans les conditions de Gauss ? Quel type d’aberrations doit-
on corriger ? Quel est l’ordre de grandeur du diamètre de la monture de l’objectif ?

\item Déterminer la position et la taille du cercle oculaire défini comme l’image de la monture
de l’objectif à travers l’oculaire. Quel est l’intérêt de placer l’œil dans le plan du cercle oculaire ?
\end{quest}
\end{exo}
%%%%%%%%%%%%%%%%%%%%%%%%%%%%%%%%%%%%%%%%%%%%%%%%%%%%%


%%%%%%%%%%%%%%%%%%%%%%%%%%%%%%%%%%%%%%%%%%%%%%%%%%%%%
\begin{exo}[PCSC][1][devoir]{Autour du microscope}
Un microscope optique est modélisable par deux lentilles minces convergentes $\mc{L}_1$ (objectif de distance focale
$f_1'$) et $\mc{L}_2$ (oculaire de distance focale $f_2'$). Il est réglé pour donner une image à l’infini d’un
objet réel AB perpendiculaire à l’axe optique (A étant sur l’axe optique) légèrement en avant
du foyer objet de l’objectif. Cette image est observée par un œil emmétrope placé au
voisinage du foyer image de l’oculaire. On notera A'B' l’image intermédiaire.
Le microscope porte les indications suivantes :
	\begin{liste}
	\item \fbox{$\times 40$} pour l’objectif, i.e. la valeur absolue du grandissement de l’objet AB par l’objectif est de 
	\begin{equation}
		\gamma_1 = 40~;
	\end{equation}
	\item \fbox{$\times 10$} pour l’oculaire, ce qui signifie que le grossissement commercial vaut 
\begin{equation}
	G_2 = 10,
\end{equation}
Le grossissement commercial est le rapport entre l’angle 
sous lequel on voit l’image à l’infini d’un objet à travers l’oculaire seul et l’angle sous lequel
on voit ce même objet à l’œil nu lorsqu’il est situé à la dimension minimale de vision distincte $d=25\,$cm ;
\item  \fbox{$\Delta = 16\,$cm} pour l’intervalle optique, i.e. la distance 
\begin{equation}
	\Delta = \overline{F_1'F_2},
\end{equation}
entre le foyer image
$F_1'$ de l’objectif et le
foyer objet $F_2$ de l’oculaire.
\item \fbox{$\omega_0 = 0,\!65$} pour l’ouverture numérique, i.e. la valeur de 
\begin{equation}
	\omega_0 = n \sin u,
\end{equation}
 avec $n$ l'indice de réfraction du milieu dans lequel se trouve l’objectif et $u$ l’angle maximum des rayons issus de A arrivant sur l’objectif.
	\end{liste}

\begin{quest}
	\item\label{quest:schema_micro} Faire un schéma du dispositif (sans respecter l’échelle) et tracer la marche de trois rayons
lumineux issus du point B de l’objet AB, l’un émis parallèlement à l’axe optique, un autre
passant par le foyer objet de l’objectif, et au autre passant par le centre de la lentille $L_1$.
\sol{
	\begin{fig}[label]
		\centering
		\includegraphics[width=0.8\columnwidth]{trace_micro.png}
		%\capt{}
	\end{fig}
}
	\item En utilisant le grossissement $G_2$ de l'oculaire, déterminer la distance focale $f_2'$ de l’oculaire.

	\sol{
		On travaille dans les \udl{conditions de Gauss} : tous les angles sont petits. Considérons un objet $\overline{A_1B_1}$ vu par l'oculaire seul.  S'appuyer sur un schéma séparé pour les calculs ci-dessous.
		
		\begin{enumerate}
			\item 
			En présence de l'oculaire, l'angle sous lequel est vu l'objet vaut 
			$$
			\alpha' \simeq \tan\alpha' = \frac{\overline{A_1B_1}}{f_2'}.
			$$
			\item En l'absence de l'oculaire, l'angle sous lequel est vu l'objet placé au \emph{punctum proximum}  vaut 
			$$
			\alpha_{pp} \simeq \tan \alpha_{pp} = \frac{\overline{A_1B_1}}{d}.
			$$
		\end{enumerate}
		Par définition du \udl{grossissement commercial}, on a donc
		$$
		G_2 = \frac{\alpha'}{\alpha_{pp}} = \frac{d}{f_2'}.
		$$
		On en déduit 
		\eqbox{f_2'=\frac{d}{G_2}=\SI{2,5}{cm}.}
	}

	\item En s'aidant de la formule de conjugaison de Newton, et en utilisant le grandissement $\gamma_1$ de l’objectif, déterminer la distance focale $f_1'$ de l’objectif.

\sol{
%\includegraphics[scale=0.2]{g1.png}
On peut utiliser une formule de conjugaison. Le plus simple est d'utiliser la formule de \udl{conjugaison de Newton} :
$$
\overline{F_1'A_1} \cdot \overline{F_1A} =  - f_1'^2.
$$
Avec $A_1=F_2$, on a donc 
$$
\Delta \cdot \overline{F_1A} = - f_1'^2.
$$
 Par définition du \udl{grandissement} de l'objectif, et d'après le \udl{théorème de Thalès} :
 $$
	\gamma_1 = \frac{\overline{O_1H}}{\overline{AB}} = \frac{\overline{F_1O_1}}{\overline{F_1A}} = \frac{f_1'}{\overline{F_1A}},
 $$ d'où $\gamma_1 = \Delta/f_1'$ et donc
\eqbox{
f_1'=\frac{\Delta}{\gamma_1}=\SI{4}{mm}.
}
}


\item Calculer la distance $\overline{O_1A}$ permettant de positionner l’objet.

\sol{
%\includegraphics[scale=0.2]{O1A.png}
 À partir du résultat ci-dessus on trouve
 \eqbox{
	 \overline{F_1 A} =-\frac{f_1'^2}{\Delta}=-\SI{0.1}{mm},
	 }
	  puis d'après la \udl{relation de Chasles}, $$\overline{O_1 A} = \overline{O_1 F_1} + \overline{F_1A} = f_1' + \overline{F_1A},$$ soit
 \eqbox{
\overline{O_1 A}  =-\SI{4.1}{mm}.
 }
}

\item Calculer le grossissement commercial $G$ du microscope.
\sol{
	On raisonne par des considérations géométriques sur le schéma de la \Quest{schema_micro}.
	\begin{enumerate}
		\item 
		L'angle de l'image en sortie du microscope vaut, toujours dans les conditions de Gauss,
		$$\alpha' \simeq  \frac{\overline{A_1B_1}}{\overline{F_2O_2}} = \frac{\overline{O_1H}}{f_2'} = \frac{\overline{AB}}{\overline{F_1A}}\cdot \frac{f_1'}{f_2'}.$$
		\item L'angle de vision au \emph{punctum proximum} vaut $$\alpha \simeq \frac{\overline{AB}}{d}.$$ 
	\end{enumerate}
	On en déduit le \udl{grossissement commercial} du microscope, qui vaut par définition
$$
G = \frac{\alpha'}{\alpha}, 
$$
soit
\eqbox{
G = \frac{d}{\overline{F_1A}}\cdot \frac{f_1'}{f_2'} = \SI{400}{}.
}
%\includegraphics[scale=0.2]{a.png}\includegraphics[scale=0.2]{ap.png}
%\includegraphics[scale=0.2]{AF2.png}\includegraphics[scale=0.2]{G.png}
}
% \item Déterminer la latitude de mise au point, à savoir la variation de la distance $\overline{O_1A}$ permettant de passer d'une observation à l'infini (ce qui a été supposé jusqu'à présent) à une observation à la distance minimale que l'œil peut accomoder, à savoir $d=25\,$cm. L'image A'B' se trouve alors à une distance $d$ à gauche de l'oculaire (elle est virtuelle).
%
\item Calculer l’angle $u$ intervenant dans l’ouverture numérique pour un objectif placé dans
l’air. Le microscope est-il utilisé dans les conditions de Gauss ? Quel type d’aberrations doit-on corriger ? Quel est l’ordre de grandeur du diamètre de la monture de l’objectif ?
\sol{
	\begin{enumerate}
		\item 
		On trouve, avec $n\simeq1$ pour l'air,
		\eqbox{
			u = \arcsin(\omega_0/n) = \SI{40}{^\circ}.
		}
		\item Le microscope n'est \udl{pas utilisé} dans les conditions de Gauss, car $u$ n'est pas petit !
		\item Lorsque les conditions de Gauss ne sont pas respectées, on doit corriger les \udl{aberrations géométriques} : aberration sphérique, coma, astigmatisme, courbure de champ et distorsion.
		\item Le diamètre de la monture de l'objectif est de l'ordre de
		\eqbox{
			D_1 \simeq 2 f_1' \tan u \simeq \SI{5.8}{mm}.
		}
	\end{enumerate}
}
%
% \item Déterminer la position et la taille du cercle oculaire défini comme l’image de la monture
% de l’objectif à travers l’oculaire. Quel est l’intérêt de placer l’œil dans le plan du cercle oculaire ?
\end{quest}
\end{exo}
%%%%%%%%%%%%%%%%%%%%%%%%%%%%%%%%%%%%%%%%%%%%%%%%%%%%%


%%%%%%%%%%%%%%%%%%%%%%%%%%%%%%%%%%%%%%%%%%%%%%%%%%
\begin{exo}[PCSI][1][colle]{Lunette astronomique}
On associe deux lentilles convergentes de distances focales $f^{'}_1$ et $f^{'}_2$.
On les dispose de manière à ce que le foyer image (principal) de $\mc{L}_1$  coïncide avec le foyer objet de $\mc{L}_2$.
On s'intéresse ici à un objet situé à l'infini et ayant dont les rayons arrivent inclinés d'un angle $\alpha$ avec l'axe optique.
\begin{quest}
\item Représenter le système ainsi réalisé, en indiquant le centre de chaque lentille ainsi que les foyers images et objets de chaque lentille. On prendra $f^{'}_1= 1$ cm et $f^{'}_2=5$ cm.
\item Tracer l'allure, à travers ce système optique, du chemin suivi par un rayon venant de l'infini en formant un angle $\alpha$ avec l'axe optique.
\item Où se situe l'image d'un point situé à l'infini sur l'axe optique par ce système ?
\item Ce rayon sort de la lunette avec un angle $\beta$. Calculer le grossissement de la lunette $G=\frac{\beta}{\alpha}$.
\end{quest}
\end{exo}
%%%%%%%%%%%%%%%%%%%%%%%%%%%%%%%%%%%%%%%%%%%%%%%%%%%%%%


%%%%%%%%%%%%%%%%%%%%%%%%%%%%%%%%%%%%%%%%%%%%%%%%%%%%%%%%%%%
\begin{exo}[PCSI]{\'Etude d'un viseur}
Un viseur est constitué d'une lentille $(L_{1},f'_{1}=10$ cm) convergente formant l'objectif et d'un oculaire $(L_{2},f'_{2}=2$ cm). La distance $D$ entre les lentilles
est réglable.

\begin{quest}
 \item Comment régler $D$ pour que le système soit afocal ? Représenter la marche d'un rayon.
 \item Quelle est le grandissement angulaire $G$ dans ces conditions ?
 \item On régle le viseur pour qu'un observateur voit nettement ; sans accomoder, u objet situé à 20 cm de l'objectif. Quelle est la nouvelle valeur de $D$ ?
\end{quest}
\end{exo}
%%%%%%%%%%%%%%%%%%%%%%%%%%%%%%%%%%%%%%%%%%%%%%%%%%%%%%%%%%%%%%%

%%%%%%%%%%%%%%%%%%%%%%%%%%%%%%%%%%%%%%%%%%%%%%%%%%
\begin{exo}[PCSI]{Élargisseur de faisceau}
On associe deux lentilles convergentes de distances focales $f^{'}_1=5$ cm et $f^{'}_2=10$ cm
On les dispose de manière à ce que le foyer image (principal) de $L_1$  coïncide avec le foyer objet de $L_2$ .
\begin{quest}
\item Représenter le système ainsi réalisé, en indiquant le centre de chaque lentille ainsi que les foyers images et objets de chaque lentille.
\item Tracer l'allure, à travers ce système optique, de deux rayons arrivant parallèles à l'axe optique.
\item Où se situe l'image d'un point situé à l'infini sur l'axe optique par ce système ?
\item Ce système est en fait un élargisseur de faisceau. On envoie alors un faisceau lumineux parallèle de diamètre $d$ dans le système optique. Donner le rapport de la largeur de sortie du faisceau sur celle d'entrée en fonction de $f^{'}_1$  et $f^{'}_2$.
\end{quest}
\end{exo}
%%%%%%%%%%%%%%%%%%%%%%%%%%%%%%%%%%%%%%%%%%%%%%%%%



%%%%%%%%%%%%%%%%%%%%%%%%%%%%%%%%%%%%%%%%%%%%%%%%%%%%%%%%%%%
\begin{exo}[PCSI]{Loupe}
Dans cet exercice, on s'intéresse à un système centré, l'axe optique étant orienté positivement dans le sens de propagation de la lumière. Un observateur émmétrope peut voir distinctement des objets situés à une distance comprise entre $d_m$ et l'infini.
\begin{quest}
\item Un observateur regarde à l'œil nu un tout petit objet plan, qu'on assimile à un segment $AB$ de longueur $\ell$, orthogonal à l'axe optique. Proposer un modèle optique de l'œil. Définir les punctum proximum et punctum remotum. Déterminer l'angle maximal $\alpha_m$ sous lequel peut être vu l'objet par un œil émmétrope.

\sol{
$\alpha_m \simeq \tan \alpha_m = \ell/d_m$.
}

\item L'observateur regarde $AB$ à travers une lentille mince convergente, de distance focale $f'$ et de centre $\rm O$ (loupe). Son œil est situé à une distance $a<d_m$ de la loupe.
\begin{enumerate}
\item Déterminer les positions de l'objet rendant possible l'observation d'une image nette. Faire une construction géométrique de l'image. L'image est-elle droite ou renversée ?
\sol{
 En notant $x' = \overline{OA'}$ avec $A'$ l'image de $A$ et $x = \overline{OA}$, l'observation n'est nette que si $a-x'$ (distance entre l'image et l'œil) est supérieur à $d_m$. La relation de conjugaison peut s'écrire :
\begin{equation*}
x' = \frac{xf'}{x+f'}
\end{equation*}
Donc la condition est :
\begin{equation*}
a - \frac{xf'}{x+f'} \geq d_m \qquad \Rightarrow \qquad a(x+f') -xf' \geq d_m(x+f')
\end{equation*}
à condition que $x+f' >0$ sinon l'égalité change de sens. Dans le cas $x+f'<0$ l'objet est avant le foyer objet, l'image est renversée, ce qui ne correspond pas à ce qu'on attend d'une loupe et ce qui conduirait de plus à une incompatibilité dans les inégalités. On obtient finalement :
\begin{equation*}
-f' \leq x \leq - \frac{f'(d_m-a)}{d_m-a+f'}
\end{equation*}
}
\item Pour quelle position l'observation se fait-elle sans accommodation de la part de l'œil ? Exprimer l'angle $\alpha$ sous lequel est vu l'objet dans ce cas là.

Application numérique : déterminer le grossissement commercial de la loupe $G = \alpha/\alpha_m$. On donne $d_m = 25$ cm et $f' = 50$ mm.

\sol{
L'observation se faire sans fatigue lorsque l'image est à l'infini, donc lorsque l'objet est dans le plan focal objet de la lentille : $\alpha \simeq \tan \alpha =\ell/f'$ soit :
\begin{equation*}
G = \frac{\alpha}{\alpha_m} = \frac{d_m}{f'} = 5
\end{equation*}
}

\end{enumerate}
\end{quest}

\end{exo}
%%%%%%%%%%%%%%%%%%%%%%%%%%%%%%%%%%%%%%%%%%%%%%%%%%%%%%%%%


%%%%%%%%%%%%%%%%%%%%%%%%%%%%%%%%%%%%%%%%%%%%%%%%%%%%%%%%%%%%%%%%%%%%%%%%%
\begin{exo}[PCSI]{\'Etude d'un rétroprojecteur}
Un spectroscope est un appareil destiné à étudier le spectre d'une source lumineuse. Un collimateur permet de réaliser un faisceau de rayons parallèles qui va éclairer un prisme. Un viseur permet ensuite d'étudier la lumière ayant traversé le prime.

Le prisme utilisé est caractérisé par un indice $n$ qui dépend de la longueur d'onde. Sa section droite est un triangle d'angle au sommet $A$. Le prisme est placé dans l'air dont l'indice sera pris égal à 1. Un rayon incident, dans le plan de la section droite, rencontre la face d'entrée au point $I$ sous l'angle d'incidence $i$ et l'émergeant associé ressort par l'autre face au point $J$ sous l'angle $i'$. L'angle $r$ (respectivement $r'$) est l'angle que fait le rayon $IJ$ dans le prisme avec la normale à la face d'entrée (respectivement normale à la face de sortie). Les angles ne sont pas orientés.

On suppose jusqu'à la question 4 la lumière monochromatique et l'indice du prisme sera égal à $n$.
\begin{quest}
\item En appliquant les lois de Descartes, écrire la relation entre $i$ et $r$ d'une part, et $i'$ et $r'$ d'autre part. On suppose que le prisme permet d'existence du rayon émergeant et on néglige, dans la suite, tout réflexion. Trouver une relation entre $r$, $r'$ et $A$.

\sol{
Les démonstrations sont celles du cours, il faut écrire :
\begin{equation*}
\sin i = n \sin r \qquad n \sin r' = \sin i' \qquad A = r + r'
\end{equation*}
}

\item \' Etablir la relation
\begin{equation*}
D = i + i' - A
\end{equation*}
En appliquant le principe du retour inverse de la lumière, montrer que, pour une valeur de $D$ possible donnée, il existe deux couples de solutions ($i$, $i'$). En déduire l'égalité de $i$ et de $i'$ lorsque $D$ passe par un minimum (supposé unique).

\sol{
C'est également du cours. D'après le principe retour inverse de la lumière si $i$ en entrée et $i'$ en sortie donnent une déviation $D$, alors $i'$ en entrée conduira à $i$ en sortie et donc aussi une déviation $D$. Puisque le minimum est unique, il n'existe qu'un incidence pour laquelle  il est obtenu, donc nécessairement on a $i = i'$.
}

\item Déterminer la valeur $i_0$ de $i$ correspondant au minimum de déviation en fonction de $n$ et de $A$. \'Etablir une relation entre $n$, l'angle $A$ et la déviation minimale $D_m$.

\sol{
 On a $\sin i_0 = n \sin r_0$ et d'autre part $A = 2 r_0$, donc on a :
\begin{equation*}
\sin i_0 = n \sin \left(\frac{A}{2}\right)
\end{equation*}
soit :
\begin{equation*}
i_0 = \mathrm{arcsin}\left(n \sin \left(\frac{A}{2}\right)\right)
\end{equation*}
Or on a $D_m = 2i_0 - A$, donc :
\begin{equation*}
\sin\left(\frac{D_m+A}{2}\right) = n \sin \left(\frac{A}{2}\right)
\end{equation*}
}

\item En déduire que l'indice $n$, les angles $A$ et $D_m$ vérifient une relation du type :
\begin{equation*}
n = \frac{f\left(\frac{A+D_m}{2}\right)}{f\left(\frac{A}{2}\right)}
\end{equation*}
où $f$ est une fonction que l'on précisera.

\sol{
 D'après la question précédente :
\begin{equation*}
n = \frac{\sin\left(\frac{A+D_m}{2}\right)}{\sin\left(\frac{A}{2}\right)}
\end{equation*}
La fonction $f$ cherchée est donc la fonction sinus.
}

\item On éclaire le prisme avec une lampe à vapeur de mercure, pour laquelle on a mesuré $D_m$ pour différentes longueurs d'onde et obtenu les valeurs de $n$ correspondantes, reprise dans la table ci-dessous.

\begin{center}
\begin{tabular}{|c|c|c|c|c|c|}
\hline
$\lambda$ ($\mu$m) & 0.4047 & 0.4358 & 0.4916 & 0.5461 & 0.5770 \\
\hline
$n$ & 1.803 & 1.791 & 1.774 & 1.762 & 1.757 \\
\hline
$1/\lambda^2$ ($\mu$m$^{-2}$) & 6.11 & 5.27 & 4.14 & 3.35 & 3.00 \\
\hline
\end{tabular}
\end{center}

Montrer que $n$ peut se mettre sous la forme
\begin{equation*}
n = a + \frac{b}{\lambda^2}
\end{equation*}
où $a$ et $b$ sont des constantes. Est-ce que le prisme est dispersif et pourquoi ?

\sol{
Si on trace $n$ en fonction de $1/\lambda^2$ alors les points s'alignent sur une droite et on trouve :
\begin{equation*}
a = 1.713 \qquad b = 0.01479\ \mu\mathrm{m}^2
\end{equation*}
Le milieu est dispersif car l'indice dépend de la longueur d'onde.
}

\item Pour une lampe à vapeur de cadmium, on mesure un indice égal à 1.777. En déduire la longueur d'onde et donner la couleur correspondante.

\sol{
\'A partir de $n$, la longueur d'onde est obtenue par :
\begin{equation*}
\lambda = \sqrt{\frac{b}{n-a}} = 480.7\ \mathrm{nm}
\end{equation*}
La raie du cadmium étudiée est donc bleue.
}

\end{quest}
\end{exo}
%%%%%%%%%%%%%%%%%%%%%%%%%%%%%%%%%%%%%%%%%%%%%%%%%%%%%%%%%%%%%%%%%%%


%%%%%%%%%%%%%%%%%%%%%%%%%%%%%%%%%%%%%%%%%%%%%%%%%%%%%%%%%%%%%%%%%%
\begin{exo}[PCSI]{Doubleur de focale}

\'A l'aide d'un objectif (assimilable à une lentille mince convergente $L$ de distance focale $f'$, on désire former l'image d'un objet réel $AB$ sur une plaque photographique située à une distance $D$ de l'objet.
\begin{quest}
\item Entre quelles valeurs peut varier $D$ pour qu'il en soit ainsi ? Préciser la distance $d$ séparant la lentille de la plaque photographique (on se limitera à $d\leq \frac{D}{2}$.

\sol{
La situation est la situation du cours où on forme une image réelle à partir d'un objet réel. On a alors la relation de conjugaison, avec $O$ le centre de la lentille, $A$ le point de l'objet sur l'axe et $A'$ le point de l'image sur l'axe :
\begin{equation*}
\frac{1}{\overline{OA'}} - \frac{1}{\overline{OA}} = \frac{1}{f'}
\end{equation*}
Or on a $D = \overline{AA'} = \overline{AO}+\overline{OA'}$ et $\overline{OA'} = d$ donc :
\begin{equation*}
\frac{1}{d} + \frac{1}{D-d} = \frac{1}{f'}
\end{equation*}
Soit :
\begin{equation*}
\frac{1}{D-d} = \frac{d-f'}{df'}
\end{equation*}
Donc :
\begin{equation*}
D = d + \frac{df'}{d-f'} = \frac{d^2}{d-f'}
\end{equation*}
On peut alors tracer le grapahe de $D(d)$ pour $d \geq f'$. Pour $d\rightarrow \infty$ la courbe tend vers l'asymptote $D = d+f'$. La courbe tend vers l'infini lorsque $d \rightarrow f'$ par valeurs supérieures. La courbe admet un minimum en $d_m = 2f'$ qui vaut
\begin{equation*}
D_m = \frac{d_m^2}{d_m-f'} = 4f'
\end{equation*}
Ainsi les solutions cherchées (avec $d\leq \frac{D}{2}$) vérifient donc : $f' \leq d \leq 2f'$ et donc :
\begin{equation*}
4f' \leq D < \infty
\end{equation*}
}
\item La lentille convergente $L$ restant fixe par rapport à l'objet $AB$, on recule la plaque photographique, et on intercale entre celle-ci et $L$ une lentille divergente $L_1$ de distance focale $f_1' = -6$ cm. Comment doit-on placer $L_1$ pour que le grandissement linéaire du système soit multiplié par deux ? De combien a-t-il fallu déplacer la plaque photographique ?

\sol{
 La lentille divergente $L_1$ doit transformer l'image précédente $A'B'$ en une nouvelle image $A_1B_1$ sur la plaque photographique telle que :
\begin{equation*}
|\gamma_1| = \left|\frac{\overline{A_1B_1}}{\overline{A'B'}}\right| = 2
\end{equation*}
et il faut que $A_1B_2$ soit réel.

En notant $O_1$ le centre de la lentille, on a la relation de conjugaison :
\begin{equation*}
\frac{1}{\overline{O_1A_1}} - \frac{1}{\overline{O_1A'}} = \frac{1}{f_1'} = - \frac{1}{|f_1'|}
\end{equation*}
Donc :
\begin{equation*}
\frac{1}{\overline{O_1A'}} = \frac{1}{\overline{O_1A_1}} + \frac{1}{|f_1'|} >0
\end{equation*}
Donc $\overline{O_1A'}>0$ montre que $A'B'$ joue le rôle d'objet virtuel pour $L_1$, ce qui est cohérent avec l'énoncé : on a place $L_1$ entre $L$ et la position initiale de la plaque photographique...

Le grandissement $\gamma_1$ est donc positif, et :
\begin{equation*}
\gamma_1 = + 2 = \frac{\overline{O_1A_1}}{\overline{O_1A'}} \qquad \Rightarrow \qquad \overline{O_1A_1} = 2 \overline{O_1A'}
\end{equation*}
La relation de conjugaison donne donc :
\begin{equation*}
\overline{O_1A'} = \frac{1}{2}|f_1'| \qquad \overline{O_1A_1} = |f_1'|
\end{equation*}
La plaque photographique doit se trouver dans le plan focal objet de la lentille divergente $L_1$, celle-ci étant placé à une distance $\frac{1}{2}|f_1'| = 3$ cm en avant de la position initiale de la plaque photographique.

Il faut maintenant reculer la plaque photographique d'une distance :
\begin{equation*}
\overline{O_1A_1} - \overline{O_1A_1} = \frac{1}{2}|f_1'| = 3\ \mathrm{cm}
\end{equation*}
}
\end{quest}
\end{exo}
%%%%%%%%%%%%%%%%%%%%%%%%%%%%%%%%%%%%%%%%%%%%%%%
